% 题证摘录；来源与整理边界见相同编号分题文件。

\paragraph{Weierstrass division: the existence construction used below.}
\begin{thm}[Weierstrass division theorem, original statement]
Suppose $f$ is holomorphic near the origin, and suppose $P$ is a Weierstrass polynomial of degree $k\geq1$ in $z_n$. Then there exists a neighborhood $V$ of the origin and unique $q,r\in\sO(V)$, where $r$ is a polynomial in $z_n$ of degree less than $k$, and on $V$,
\[f=qP+r.\]
\end{thm}
Note that $r$ need not be a Weierstrass polynomial; it need not be monic nor do the coefficients need to vanish at the origin. It is simply a polynomial in $z_n$ with coefficients that are holomorphic functions of the first $n-1$ variables.
\begin{proof}
Uniqueness is left as an exercise. There exists a connected neighborhood $V=V'\times D$ of the origin, where $D$ is a disc, $f$ and $P$ are continuous on $V'\times\widebar D$, and $P$ is not zero on $V'\times\partial D$. Let
\[
q(z',z_n)=\frac1{2\pi i}\int_{\partial D}\frac{f(z',\zeta)}{P(z',\zeta)(\zeta-z_n)}\,d\zeta.
\]
As $P$ is not zero on $V'\times\partial D$, the function $q$ is holomorphic in $V$ (differentiate under the integral). If $P$ did divide $f$, then $q$ would really be $\frac fP$. But if $P$ does not divide $f$, then the Cauchy integral formula does not apply and $q$ is not equal to $\frac fP$. Interestingly, the expression does give the quotient in the division with remainder. Write $f$ using the Cauchy integral formula in $z_n$ and subtract $qP$ to obtain $r$:
\begin{align*}
r(z',z_n)&=f(z',z_n)-q(z',z_n)P(z',z_n)\\
&=\frac1{2\pi i}\int_{\partial D}\frac{f(z',\zeta)P(z',\zeta)-f(z',\zeta)P(z',z_n)}{P(z',\zeta)(\zeta-z_n)}\,d\zeta.
\end{align*}
We need to show $r$ is a polynomial in $z_n$ of degree less than $k$. In the expression inside the integral, the numerator is of the form $\sum_{\ell=1}^k h_\ell(z',\zeta)(\zeta^\ell-z_n^\ell)$ and is therefore divisible by $(\zeta-z_n)$. The numerator is a polynomial of degree $k$ in $z_n$. After dividing by $(\zeta-z_n)$, the integrand becomes a polynomial in $z_n$ of degree $k-1$. Use linearity of the integral to integrate the coefficients of the polynomial. Each coefficient is a holomorphic function in $V'$ and the proof is finished. Some coefficients may have integrated to zero, so we can only say that $r$ is a polynomial of degree $k-1$ or less.
\end{proof}

\paragraph{Noetherian property (the counted problem).}
A commutative ring $R$ is Noetherian if every ideal in $R$ is finitely generated. That is, for every ideal $I\subset R$ there exist finitely many generators $f_1,\ldots,f_k\in I$: Every $g\in I$ can be written as $g=c_1f_1+\cdots+c_kf_k$, for some $c_1,\ldots,c_k\in R$.
\begin{thm}
$\sO_p$ is Noetherian.
\end{thm}
\begin{proof}
Without loss of generality, $p=0$. The proof is by induction on dimension. By \exerciseref{exercise:onedimOisPID}, ${}_1\sO_0$ is Noetherian. By \exerciseref{exercise:biholringiso}, we are allowed a biholomorphic change of coordinates near the origin. For induction, suppose ${}_{n-1}\sO_0$ is Noetherian and let $I\subset{}_n\sO_0$ be an ideal. If $I=\{0\}$ or $I={}_n\sO_0$, then the assertion is obvious. Therefore, assume that all elements of $I$ vanish at the origin ($I\neq{}_n\sO_0$), and that there exist elements that are not identically zero ($I\neq\{0\}$). Let $g$ be such an element. After perhaps a linear change of coordinates, assume $g$ is a Weierstrass polynomial in $z_n$ by the preparation theorem. The ring ${}_{n-1}\sO_0[z_n]$ is a subring of ${}_n\sO_0$. The set $J={}_{n-1}\sO_0[z_n]\cap I$ is an ideal in the ring ${}_{n-1}\sO_0[z_n]$. By the Hilbert basis theorem (see \thmref{thm:hilbertbasis} in the appendix for a proof), as ${}_{n-1}\sO_0$ is Noetherian, the ring ${}_{n-1}\sO_0[z_n]$ is also Noetherian. Thus $J$ has finitely many generators, that is, $J=(h_1,\ldots,h_k)$ in the ring ${}_{n-1}\sO_0[z_n]$. By the division theorem, every $f\in I$ is of the form $f=qg+r$, where $r\in{}_{n-1}\sO_0[z_n]$ and $q\in{}_n\sO_0$. As $f$ and $g$ are in $I$, so is $r$. As $g$ and $r$ are in ${}_{n-1}\sO_0[z_n]$, they are both in $J$. Write $g=c_1h_1+\cdots+c_kh_k$ and $r=d_1h_1+\cdots+d_kh_k$. Then
\[f=(qc_1+d_1)h_1+\cdots+(qc_k+d_k)h_k.\]
So $h_1,\ldots,h_k$ also generate $I$ in ${}_n\sO_0$.
\end{proof}
